\section{Related Work}

Change point detection (CPD) methods are usually grouped into online and
offline formulations; the offline setting used here estimates event locations
after the full series has been observed \cite{adams2007}. Surveys emphasize
that benchmark conclusions depend on detector assumptions, calibration,
post-processing, data generation, annotation, and matching rules
\cite{aminikhanghahi2017,truong2020,vandenburg2020}. Most simulations prescribe
the change structure directly, for example by placing shifts in a piecewise
signal or concatenating distinct segments
\cite{killick2012,fryzlewicz2014,gharghabi2019,ermshaus2023clasp}. Here the
labels come from a simulated stochastic trajectory: their number and spacing
are consequences of the dynamics and noise realization.

This distinction separates the benchmark from SDE-based detection models.
Recent work has used latent stochastic differential equations as a model class
for change point detection \cite{ryzhikov2023latent_sde_cpd}. Here the
stochastic dynamics are not the detector. They define the transition process on
which existing classical, window-based, and learned detectors are evaluated
under one fixed protocol.

The stochastic-dynamics motivation is rooted in noise-induced motion between
regions of state space. Classical work on random perturbations, escape
problems, and sample paths describes how fluctuations can drive transitions in
nonlinear systems \cite{kramers1940,freidlin1998,berglund2006}. Within
\textit{Fluctuation and Noise Letters}, noise-induced escape from the Lorenz
attractor has been studied as an escape problem in nonlinear dynamics
\cite{anishchenko2001fnl_escape}. Colored noise is also a natural part of this
setting rather than only a robustness stress test: FNL studies have examined
the role of colored noise in stochastic-system behavior
\cite{valenti2005fnl_colored}, and recent work has used deep learning to
characterize colored-noise time-series patterns
\cite{barauna2024fnl_colored_deep}. These papers motivate treating the
colored-noise diagnostic in this article as a boundary on neural claims, not as
a generic machine-learning add-on.

The compared methods cover standard CPD families: CUSUM-type, rank-based,
homogeneity, and structural-break tests
\cite{page1954,pettitt1979,alexandersson1986,buishand1982,chow1960};
segmentation and local-window methods
\cite{killick2012,harchaoui2009,liu2013}; and learned detectors trained from
labeled windows \cite{li2024deepcpd}. Because learned detectors use generated
labels, the benchmark keeps training, calibration, and test trajectories
independent.
